% !TEX root = ../main.tex

\chapter{Marco Teórico}

Este capítulo presenta los fundamentos conceptuales necesarios para comprender tanto el problema abordado en esta memoria como la solución desarrollada. En primer lugar, se introducen las características de \textsf{Haskell} necesarias para interpretar los ejemplos y las firmas de tipos utilizadas a lo largo del informe. Sobre esta base, se estudian sus principales abstracciones funcionales: functores, aplicativos y mónadas. Estas clases forman una jerarquía que permite estructurar cómputos dentro de un contexto y se distinguen por el grado de dependencia que pueden expresar entre las distintas partes de un programa.

A continuación, se profundiza en la \textit{do-notation}, explicando su relación con la composición monádica y cómo sus statements de enlace se traducen, en su forma convencional, mediante el operador \textit{bind} \verb|(>>=)|. Esto es fundamental para entender por qué \texttt{ApplicativeDo} puede reemplazar parte de una expresión monádica por una expresión aplicativa \cite{MarlowPJKM2016,GHCApplicativeDo}.

Sobre esta base, se introduce la noción de mónada conmutativa, cuya propiedad relevante permite intercambiar ciertos cómputos independientes sin modificar la semántica del programa \cite{Kock2012,Jacobs2018}. Posteriormente, se presentan las mónadas de probabilidad como un caso canónico de esta propiedad y como el principal dominio de aplicación de la memoria.

Finalmente, se estudia la extensión \texttt{ApplicativeDo}, describiendo la transformación que realiza sobre bloques \texttt{do}, la forma en que identifica dependencias entre statements y la generación de expresiones aplicativas que pueden exponer una mayor estructura paralelizable. Este recorrido permite establecer los conceptos necesarios para analizar posteriormente las limitaciones del mecanismo actual y la extensión experimental propuesta.

% !TEX root = ../../main.tex

\section{Fundamentos de Haskell}
\label{sec:fundamentos-haskell}

\textsf{Haskell} es un lenguaje de programación funcional puro, de propósito general y con tipado estático. En este lenguaje, las funciones desempeñan un papel central y son tratadas como valores de primera clase: pueden recibirse como argumentos, retornarse como resultados y almacenarse en estructuras de datos. Esta característica permite que los programas no solo operen con valores como números o listas, sino también con las funciones que transforman dichos valores \cite[capítulo~1 y sección~6.1.6]{Haskell2010}.

La pureza implica que la evaluación de una función depende de sus argumentos y no produce efectos secundarios implícitos. Los efectos se representan de manera explícita mediante tipos y abstracciones que determinan cómo se componen los cómputos. Esta separación será relevante al introducir mónadas y analizar el orden de los efectos expresados mediante bloques \texttt{do} \cite[capítulo~7]{Haskell2010} \cite{Wadler1995}.

Además, las funciones de \textsf{Haskell} se encuentran currificadas. Esto significa que una función que conceptualmente recibe varios argumentos se representa como una sucesión de funciones que reciben un argumento a la vez. Por esta razón, las flechas de una firma de tipos se asocian hacia la derecha; por ejemplo, el tipo \verb|a -> b -> c| se interpreta como \verb|a -> (b -> c)|. Una función de este tipo recibe un valor de tipo \verb|a| y retorna una nueva función de tipo \verb|b -> c| \cite[secciones~3.3 y~4.1.2]{Haskell2010}.

Este comportamiento puede observarse en el operador binario \verb|(+)|, cuya firma de tipos se muestra en el Código~\ref{lst:sum-type-signature} \cite[sección~6.4 y figura~6.2]{Haskell2010}.

\begin{lstlisting}[style=haskell,caption={Firma de tipos del operador de suma \texttt{(+)}.},label={lst:sum-type-signature}]
    (+) :: Num a => a -> a -> a
\end{lstlisting}

En esta firma, \verb|a| es una variable de tipo y la restricción \verb|Num a| establece que debe elegirse un tipo perteneciente a la clase \texttt{Num}. Si se aplica \verb|(+)| únicamente al valor \verb|5|, se fija su primer argumento y se obtiene una nueva función que todavía espera el argumento restante. Esta operación se denomina \textit{aplicación parcial}. En este caso, la expresión \verb|(+) 5| produce la función unaria \verb|(5+)| de tipo \verb|Num a => a -> a| \cite[secciones~3.5, 4.1.3 y~6.4]{Haskell2010}.

Los tipos también pueden estar parametrizados. Por ejemplo, \texttt{Maybe} es un constructor de tipos que recibe un tipo \verb|a| y produce el tipo \texttt{Maybe a}. Sus constructores de datos representan dos posibilidades: \texttt{Just} contiene un valor de tipo \verb|a| y \texttt{Nothing} representa su ausencia. Esta distinción entre un constructor de tipos y los tipos concretos que produce permite definir comportamientos comunes para toda la familia \texttt{Maybe} \cite[secciones~4.1.2 y~6.1.8; capítulo~21]{Haskell2010}.

Una clase de tipos especifica un conjunto de operaciones que distintos tipos o constructores de tipos pueden implementar. Cada implementación concreta se declara mediante una instancia \cite[secciones~4.1 y~4.3]{Haskell2010}. La restricción \verb|Num a| del ejemplo anterior indica precisamente que debe existir una instancia de \texttt{Num} para el tipo elegido. De manera análoga, \texttt{Functor}, \texttt{Applicative} y \texttt{Monad} son clases de tipos con operaciones e instancias propias. \texttt{Maybe} posee instancias para las tres y, por ello, se utilizará como ejemplo común para mostrar cómo un mismo contexto puede transformar valores, combinar cómputos y expresar dependencias entre sus resultados \cite{Yorgey2009,GHCInternalBase}.

% !TEX root = ../../main.tex

\section{Abstracciones Funcionales}

Las funciones de primera clase, la currificación, la aplicación parcial y las clases de tipos proporcionan la base para las abstracciones estudiadas en esta sección. \texttt{Functor}, \texttt{Applicative} y \texttt{Monad} forman una jerarquía que se distingue por el grado de dependencia que permiten expresar entre las partes de un programa. Un functor permite transformar el resultado de un cómputo; un aplicativo permite combinar cómputos cuya estructura está determinada de antemano; y una mónada permite que un cómputo posterior dependa del valor producido por uno anterior \cite{Wadler1995,McBride2008,Yorgey2009}.

\subsection{Functores}

En \textsf{Haskell}, \texttt{Functor} es una clase de tipos cuyas instancias se definen para constructores de tipos, como las listas \verb|[]| o \verb|Maybe|. Su operación principal es \texttt{fmap}, que permite aplicar una función a los valores asociados a una estructura y obtener un resultado que conserva el mismo constructor de tipos. La definición simplificada de la clase se muestra en el Código~\ref{lst:functor-class}:

\begin{lstlisting}[style=haskell,caption={Definición simplificada de la clase \texttt{Functor}.},label={lst:functor-class}]
class Functor f where
    fmap :: (a -> b) -> f a -> f b
\end{lstlisting}

En esta firma, \verb|f| representa un constructor de tipos, mientras que \verb|a| y \verb|b| representan los tipos de los valores antes y después de aplicar la función. Así, \texttt{fmap} recibe una función de tipo \verb|a -> b| y transforma un valor de tipo \verb|f a| en uno de tipo \verb|f b|, manteniendo el contexto representado por \verb|f|.

En instancias concretas como las listas y \verb|Maybe|, esta operación puede entenderse intuitivamente mediante la metáfora de un contenedor: \texttt{fmap} transforma los valores contenidos sin modificar la estructura que los organiza. Bajo esta interpretación pedagógica, puede imaginarse que la operación accede a los valores, aplica una función sobre ellos y vuelve a representarlos dentro del mismo contexto. Sin embargo, esta metáfora no constituye la definición general de un functor.

Además, toda instancia válida de \texttt{Functor} debe respetar las leyes de identidad y composición. Estas leyes garantizan que la aplicación de \texttt{fmap} preserve de manera coherente la estructura asociada al constructor de tipos \cite{Yorgey2009}.

Los ejemplos clásicos de functores incluyen las listas y los árboles, pues \texttt{fmap} transforma sus elementos de manera uniforme mientras mantiene su estructura subyacente.

Como se mencionó anteriormente, otro ejemplo relevante es \texttt{Maybe}. El tipo \texttt{Maybe a} representa valores opcionales: un valor disponible se expresa mediante el constructor de datos \texttt{Just}, mientras que la ausencia de un valor se representa mediante \texttt{Nothing}. Al utilizar \texttt{fmap} sobre este tipo, la función recibida se aplica al valor contenido en \texttt{Just}, mientras que \texttt{Nothing} permanece sin cambios.

En la implementación de \texttt{base} incluida con GHC, este comportamiento se define mediante la instancia de \texttt{Functor} para el constructor de tipos \texttt{Maybe} mostrada en el Código~\ref{lst:functor-maybe-instance} \cite{GHCInternalBase}:

\begin{lstlisting}[style=haskell,caption={Instancia de \texttt{Functor} para \texttt{Maybe}.},label={lst:functor-maybe-instance}]
instance Functor Maybe where
    fmap _ Nothing  = Nothing
    fmap f (Just a) = Just (f a)
\end{lstlisting}

La primera ecuación establece que, cuando el valor recibido es \texttt{Nothing}, no existe ningún elemento sobre el cual aplicar la función y, por tanto, el resultado permanece como \texttt{Nothing}. El guion bajo indica que la función recibida no necesita utilizarse en este caso. En cambio, cuando el valor tiene la forma \texttt{Just a}, la segunda ecuación aplica la función \texttt{f} al elemento \texttt{a} y vuelve a encapsular el resultado mediante el constructor \texttt{Just}. Ambos comportamientos se ilustran en el Código~\ref{lst:functor-maybe-fmap-examples}:

\begin{lstlisting}[style=haskell,caption={Aplicación de \texttt{fmap} sobre valores de tipo \texttt{Maybe}.},label={lst:functor-maybe-fmap-examples}]
    > fmap (+1) (Just 2)
    Just 3
    > fmap (+1) Nothing
    Nothing
\end{lstlisting}

\subsection{Aplicativos}

En \textsf{Haskell}, \texttt{Applicative} es una clase de tipos definida para constructores de tipos que ya poseen una instancia de \texttt{Functor}. Sus operaciones principales son \texttt{pure} y \verb|(<*>)|: la primera introduce un valor en un contexto, mientras que la segunda aplica una función contextualizada a un valor dentro de ese mismo contexto. De esta forma, la clase caracteriza un estilo aplicativo de programación, menos expresivo que el monádico y, por tanto, aplicable a una mayor variedad de estructuras \cite{McBride2008}. Para ilustrar este comportamiento, se utilizará nuevamente el constructor de tipos \texttt{Maybe} y el operador binario \texttt{(+)}.

Un constructor de tipos que posee una instancia de \texttt{Applicative} se denomina aplicativo. Considerando \texttt{Maybe}, para sumar \texttt{Just 2} y \texttt{Just 3}, la función \texttt{(+)} se introduce primero en el contexto mediante \texttt{pure}. La primera aplicación de \verb|(<*>)| produce la función unaria contextualizada \texttt{Just ((+) 2)}, de tipo \verb|Maybe (Int -> Int)|. La segunda aplicación combina esta función con \texttt{Just 3} y produce el resultado \texttt{Just 5}.

Cada instancia de \texttt{Applicative} define cómo se comportan estas operaciones para un constructor de tipos concreto. En la instancia de \texttt{Maybe}, si alguno de los cómputos asociados a los sumandos produce \texttt{Nothing}, esta ausencia se propaga al resultado. De esta forma, el comportamiento de los valores opcionales queda definido una sola vez en la instancia, evitando repetir ajustes de patrones cada vez que se combinan operaciones que pueden fallar.

Formalmente, la clase \texttt{Applicative} extiende \texttt{Functor} y declara las operaciones \texttt{pure} y \verb|(<*>)| mediante la definición simplificada del Código~\ref{lst:applicative-class}:

\begin{lstlisting}[style=haskell,caption={Definición simplificada de la clase \texttt{Applicative}.},label={lst:applicative-class}]
    class Functor f => Applicative f where
        pure :: a -> f a
        <*>  :: f (a-> b) -> f a -> f b
\end{lstlisting}

En particular, la instancia \verb|Applicative Maybe| define los operadores \verb|pure| y \verb|<*>| de acuerdo con el comportamiento de los valores opcionales: permite introducir valores en el contexto y combinar operaciones que pueden fallar. Una definición de esta instancia se presenta en el Código~\ref{lst:applicative-maybe-instance}:

\begin{lstlisting}[style=haskell,caption={Instancia de \texttt{Applicative} para \texttt{Maybe}.},label={lst:applicative-maybe-instance}]
    instance Applicative Maybe where
        pure = Just
        Nothing <*> _ = Nothing
        _ <*> Nothing = Nothing
        (Just f) <*> (Just x) = Just (f x)
\end{lstlisting}

Nótese que, en caso de operar con algún fallo de cómputo o ausencia de algún valor, cuya abstracción equivale al constructor \verb|Nothing|, las siguientes aplicaciones retornarán a su vez un elemento \verb|Nothing|, propagando así el comportamiento definido por la instancia \verb|Applicative Maybe|. En caso contrario, si ambos operandos se encuentran definidos, es decir, son elementos encapsulados en el constructor \verb|Just|, se procede con la aplicación del primero sobre el segundo, para finalmente empaquetar el elemento resultante dentro del mismo constructor \verb|Just|.

Volviendo al ejemplo de aplicación de la suma \verb|(+)| sobre elementos de tipo \verb|Maybe Int|, se requiere que esta se encuentre encapsulada dentro del constructor \verb|Just|. El operador \verb|<*>| espera que su operando izquierdo sea una función dentro del contexto. Por ello, \verb|pure (+)| introduce la función de suma en \verb|Maybe|, produciendo el valor contextualizado requerido. La expresión aplicativa concreta se muestra en el Código~\ref{lst:applicative-maybe-sum-pure}:

\begin{lstlisting}[style=haskell,caption={Suma aplicativa de dos valores de tipo \texttt{Maybe Int}.},label={lst:applicative-maybe-sum-pure}]
    pure (+) <*> Just 2 <*> Just 3
\end{lstlisting}

Como \texttt{Functor} es una superclase de \texttt{Applicative}, todo constructor de tipos con una instancia de \texttt{Applicative} también dispone de \texttt{fmap}. El operador \verb|(<$>)| corresponde a su forma infija, por lo que se cumple la igualdad \verb|f <$> m = fmap f m|. Además, la coherencia entre ambas clases establece que \verb|fmap f m = pure f <*> m|. En consecuencia, ambas expresiones pueden relacionarse mediante \verb|f <$> m = fmap f m = pure f <*> m|. De este modo, el ejemplo anterior puede escribirse con la sintaxis más directa del Código~\ref{lst:applicative-maybe-sum-infix}:

\begin{lstlisting}[style=haskell,caption={Forma infija de la suma aplicativa sobre \texttt{Maybe}.},label={lst:applicative-maybe-sum-infix}]
    (+) <$> Just 2 <*> Just 3
\end{lstlisting}

Finalmente, nótese que en una expresión aplicativa, la estructura del cómputo queda fijada de antemano: cada argumento se evalúa en su propio contexto y luego se combinan sus resultados mediante la aplicación de una función pura. Esta característica hace explícita la ausencia de dependencias de datos entre los subcómputos asociados a los argumentos, pues ninguno puede usar el resultado del otro para determinar su propia forma. Esto no implica que sus efectos sean conmutativos, aunque puede habilitar una evaluación más paralelizable cuando la implementación lo permite \cite{McBride2008}.

Esta misma restricción implica que, si se quisiera que el segundo argumento dependa del valor producido por el primero, el estilo aplicativo resulta insuficiente, ya que la forma del cómputo no puede adaptarse dinámicamente a resultados intermedios. Por otro lado, la abstracción monádica proporciona el operador \textit{bind} \verb|(>>=)|, que permite construir el segundo cómputo a partir del resultado del primero. Esta capacidad de expresar dependencias entre resultados intermedios motiva la introducción de la clase \texttt{Monad} en \textsf{Haskell} \cite{Wadler1995}.

\subsection{Mónadas}

En \textsf{Haskell}, \texttt{Monad} es una clase de tipos cuyas instancias definen cómo secuenciar cómputos dentro de un mismo contexto cuando un cómputo posterior puede depender del valor producido por uno anterior \cite{Wadler1995}. Como \texttt{Applicative} es una superclase de \texttt{Monad}, todo constructor de tipos que posea una instancia de \texttt{Monad} también debe poseer una instancia de \texttt{Applicative}. El Código~\ref{lst:monad-class} muestra una definición simplificada con las operaciones relevantes:

\begin{lstlisting}[style=haskell,caption={Definición simplificada de la clase \texttt{Monad}.},label={lst:monad-class}]
class Applicative m => Monad m where
    (>>=)  :: m a -> (a -> m b) -> m b
    return :: a -> m a
\end{lstlisting}

Para interpretar la firma de \verb|(>>=)|, consideraremos un cómputo \verb|ma :: m a| y una función \verb|k :: a -> m b|. La expresión \verb|ma >>= k| compone el primer cómputo con la función k de acuerdo con el comportamiento definido por la instancia, por ende entrega su resultado de tipo \verb|a| a la función \verb|k| y obtiene un nuevo cómputo de tipo \verb|m b|. De esta manera, la función \verb|k| puede utilizar el resultado anterior para determinar la forma del cómputo siguiente.

La operación \verb|return| introduce un valor en el contexto monádico y, en la jerarquía actual de clases de tipos, coincide con \verb|pure|. Su nombre no representa una transferencia de control como en lenguajes imperativos, sino la construcción de un valor de tipo \verb|m a| a partir de uno de tipo \verb|a|.

Al encadenar varias aplicaciones de \verb|(>>=)|, las funciones entregadas al operador pueden escribirse como funciones anónimas $\lambda$. Estas permiten asignar nombres a los resultados intermedios y utilizarlos en los cómputos posteriores. Para ilustrar este comportamiento, se retomará el ejemplo de la suma con \verb|Maybe|.

En este caso, la instancia \verb|Monad Maybe| define el comportamiento del operador \verb|>>=| para los valores opcionales. Las operaciones relevantes se presentan en el Código~\ref{lst:monad-maybe-instance}:

\begin{lstlisting}[style=haskell,caption={Instancia de \texttt{Monad} para \texttt{Maybe}.},label={lst:monad-maybe-instance}]
    instance Monad Maybe where
        return x = Just x
        Nothing >>= _ = Nothing
        (Just x) >>= f = f x
\end{lstlisting}

Cuando el primer cómputo produce \verb|Nothing|, la ausencia de valor se propaga como resultado. En cambio, cuando produce un valor de la forma \verb|Just x|, el elemento \verb|x| se entrega como argumento a la función ubicada a la derecha de \verb|>>=|. En este caso, la función que se ejecuta debe respetar la firma de tipos y retornar un nuevo cómputo de tipo \verb|m b|. El Código~\ref{lst:monad-maybe-sum} presenta la aplicación de la función \verb|(+)| mediante notación monádica:

\begin{lstlisting}[style=haskell,caption={Suma monádica de dos valores de tipo \texttt{Maybe}.},label={lst:monad-maybe-sum}]
    Just 3 >>= (\x -> Just 2 >>= (\y -> return ((+) x y)))
\end{lstlisting}

Por tanto, al ejecutar el operador \verb|>>=|, se desempaqueta el número 3 y se aplica la función anónima de la derecha sobre este valor, reemplazando todas las apariciones del identificador \texttt{x} por 3. A continuación, se realiza el mismo proceso con el cómputo \verb|Just 2|, resultando finalmente en la expresión \texttt{((+) 3 2)} que será empaquetada gracias a la instrucción \texttt{return} retornando el valor contextualizado \texttt{Just 5}.

Cabe recalcar que ahora el identificador \verb|x| queda dentro del alcance de la función que construye el segundo cómputo. Por ello, la expresión asociada a \verb|Just 2| puede reemplazarse por otra que utilice \verb|x|. El Código~\ref{lst:monad-maybe-dependent-computation} muestra un programa que utiliza esta nueva expresividad de las mónadas:

\begin{lstlisting}[style=haskell,caption={Cómputo monádico con dependencia entre resultados intermedios.},label={lst:monad-maybe-dependent-computation}]
    Just 3 >>= (\x -> Just ((*) 2 x) >>= (\y -> return ((+) x y)))
\end{lstlisting}

Del programa anterior, el primer cómputo desempaqueta \verb|Just 3| y reemplaza este elemento en las apariciones del identificador \verb|x|. Luego, el cómputo posterior \verb|Just ((*) 2 x)| es capaz de considerar el valor numérico de \verb|x| retornando \verb|Just 6|, luego se desempaqueta su resultado y este reemplaza las apariciones del identificador \verb|y|. Quedando como cómputo final la operación \verb|((+) 3 6)| cuyo resultado termina siendo empaquetado por la instrucción \verb|return| como \verb|Just 9|. Este flujo de ejecución permite una mayor expresividad enriqueciendo el programa. De esta manera la notación monádica termina siendo algo característico y central en \textsf{Haskell}.

Aunque esta composición permita expresar dependencias de forma correcta y versatil, la acumulación de funciones anónimas dificulta su lectura. La \textit{do-notation}, presentada a continuación, proporciona una sintaxis más clara para escribir la misma secuencia monádica.

% !TEX root = ../../main.tex

\section{Do-Notation}
\label{sec:do-notation}

La \textit{do-notation} es una sintaxis de \textsf{Haskell} para escribir cómputos monádicos de forma secuencial. Un bloque \texttt{do} contiene statements que pueden vincular resultados intermedios y utilizarlos en los statements posteriores. Esta notación no define una forma nueva de composición, sino que proporciona una escritura más directa para expresiones construidas mediante las operaciones monádicas como puede ser el operador \textit{bind} \verb|(>>=)|, \verb|(>>)| y \texttt{fail}, además de contemplar una regla propia para declaraciones \texttt{let} \cite[sección~3.14]{Haskell2010}.

Para observar la motivación de esta sintaxis, se retomará la suma de dos valores de tipo \texttt{Maybe} presentada mediante el operador \textit{bind} en el Código~\ref{lst:monad-maybe-sum}. Aunque la expresión es compacta, la acumulación de operadores, funciones anónimas y paréntesis dificulta distinguir su estructura a simple vista.

El Código~\ref{lst:monad-maybe-sum-multiline} distribuye la misma expresión en varias líneas, sin cambiar el orden ni la agrupación de sus operaciones, para mostrar con mayor claridad el flujo de definición y composición de los cómputos:

\begin{lstlisting}[style=haskell,caption={Forma multilínea de la suma monádica sobre \texttt{Maybe}.},label={lst:monad-maybe-sum-multiline}]
    Just 3 >>= (\x ->
        Just 2 >>= (\y ->
            return ((+) x y)))
\end{lstlisting}

La indentación permite reconocer que la composición con \texttt{Just 2} se encuentra dentro de la función que introduce \texttt{x}, mientras que \texttt{return} cierra la composición utilizando los identificadores \texttt{x} e \texttt{y}. Este cambio es únicamente visual: la expresión conserva la semántica del Código~\ref{lst:monad-maybe-sum} y produce \texttt{Just 5}.

La \textit{do-notation} permite expresar este mismo flujo sin escribir de forma explícita las funciones anónimas ni las aplicaciones sucesivas de \verb|(>>=)|. En un statement de la forma \verb|p <- e|, la expresión \texttt{e} representa el cómputo y el patrón \texttt{p} vincula su resultado para que pueda utilizarse en los statements posteriores. El Código~\ref{lst:do-notation-maybe-sum} presenta la escritura equivalente:

\noindent\begin{minipage}{\textwidth}
\begin{lstlisting}[style=haskell,caption={Suma monádica sobre \texttt{Maybe} escrita mediante \textit{do-notation}.},label={lst:do-notation-maybe-sum}]
    do
        x <- Just 3
        y <- Just 2
        return ((+) x y)
\end{lstlisting}
\end{minipage}

En el Código~\ref{lst:do-notation-maybe-sum}, el primer statement corresponde al cómputo entregado al \verb|(>>=)| exterior y vincula su resultado con \texttt{x}; el segundo corresponde al \verb|(>>=)| interior y vincula su resultado con \texttt{y}; el último construye el resultado del bloque. Al aplicar las reglas convencionales de traducción de la \textit{do-notation}, este bloque recupera la composición anidada del Código~\ref{lst:monad-maybe-sum-multiline} y, en última instancia, la expresión compacta del Código~\ref{lst:monad-maybe-sum} \cite[sección~3.14]{Haskell2010}.

Aunque el bloque se escribe como una secuencia, en este ejemplo la expresión \texttt{Just 2} no utiliza el valor de \texttt{x}; por ello, entre los dos primeros statements no existe una dependencia de datos local. Ambos resultados sí son utilizados por el statement final. Esta distinción entre orden sintáctico y dependencias de datos será relevante más adelante: la notación hace más legible la composición, pero no autoriza por sí sola a intercambiar cómputos. Las condiciones que permiten transformar o reordenar partes de esta secuencia se estudiarán al introducir las mónadas conmutativas y \texttt{ApplicativeDo}.

% !TEX root = ../../main.tex

\section{Mónadas Conmutativas}
\label{sec:monadas-conmutativas}

La sección anterior mostró que la ausencia de una dependencia de datos entre dos statements permite formular dos órdenes distintos de estos mismos, pero no basta por sí sola para asegurar que el intercambio conserve el significado del programa. Las mónadas conmutativas proporcionan precisamente la propiedad semántica necesaria para justificarlo. En términos operacionales, esto significa que, dados dos cómputos monádicos independientes, ejecutar el primero y luego el segundo es semánticamente equivalente a ejecutarlos en el orden contrario, siempre que se conserven las asociaciones entre sus resultados.

Kock formaliza esta propiedad mediante dos mapas de Fubini, correspondientes a las dos formas de combinar un par de cómputos al efectuar primero uno u otro. Una mónada es conmutativa cuando ambos mapas coinciden; en ese caso, el orden elegido para realizar la combinación no modifica su significado \cite[secciones~5 y~9]{Kock2012}.

La mónada \texttt{Maybe} permite ilustrar esta propiedad mediante el mismo ejemplo de suma utilizado en las secciones anteriores. El Código~\ref{lst:commutative-maybe-sum} presenta el bloque original y una variante en la que se intercambian los dos primeros statements:

\noindent\begin{minipage}{\textwidth}
\begin{lstlisting}[style=haskell,caption={Suma monádica sobre \texttt{Maybe} en el orden original y según un reordenamiento válido.},label={lst:commutative-maybe-sum}]
    sumaOriginal = do                            sumaReordenada = do  
        x <- Just 3                                  y <- Just 2
        y <- Just 2               <==>               x <- Just 3
        return ((+) x y)                             return ((+) x y)
\end{lstlisting}
\end{minipage}

En \texttt{sumaOriginal}, el primer statement vincula el valor 3 con \texttt{x} y el segundo vincula el valor 2 con \texttt{y}. En \texttt{sumaReordenada} se invierte el orden de ejecución, pero se mantienen las mismas asociaciones: \texttt{x} continúa vinculado con 3 e \texttt{y} con 2. Como la expresión final tampoco cambia, ambos bloques producen \texttt{Just 5}.

Esta equivalencia no se debe a que la suma sea una operación conmutativa. El intercambio de statements no modifica el orden de los argumentos de \verb|((+) x y)|, sino solamente el orden en que se obtienen los valores asociados a \texttt{x} e \texttt{y}. La propiedad relevante pertenece a la mónada: bajo la semántica total de valores opcionales, si ambos cómputos producen un valor de la forma \texttt{Just}, estos valores se conservan en cualquiera de los dos órdenes; si alguno produce \texttt{Nothing}, ambos órdenes producen \texttt{Nothing}. Esta caracterización no considera valores divergentes, como \texttt{undefined}, ni efectos externos.

Por lo tanto, conviene distinguir independencia estructural de conmutatividad. La expresión \texttt{Just 2} no utiliza \texttt{x}, de modo que intercambiar ambos statements conserva su alcance y sus dependencias de datos. Sin embargo, esta ausencia de dependencia no bastaría para afirmar conmutatividad en una mónada arbitraria: dos acciones de tipo \texttt{IO} podrían no utilizar sus resultados mutuamente y aun así imprimir mensajes en órdenes diferentes. La independencia permite formular ambos órdenes, mientras que la conmutatividad garantiza que el cambio de orden no altera los efectos observables \cite[p.~94]{MarlowPJKM2016}.

La noción también aparece de manera natural en el estudio de distribuciones. Kock presenta las mónadas conmutativas como una teoría general de distribuciones \cite{Kock2012}, mientras que Jacobs identifica específicamente la mónada de distribuciones discretas finitas como una mónada conmutativa \cite[sección~3.1]{Jacobs2018}. Esta relación proporciona la base para estudiar a continuación las mónadas de probabilidad, sin implicar que toda construcción monádica probabilística sea necesariamente conmutativa.

% !TEX root = ../../main.tex

\section{Mónadas de Probabilidad}

Una mónada de probabilidad permite representar un cómputo como una distribución sobre sus posibles resultados. En vez de producir un único valor, cada operación puede introducir alternativas y asignar una probabilidad a cada una de ellas; la composición monádica combina posteriormente estas alternativas y propaga sus probabilidades. Para distribuciones discretas finitas, este enfoque permite enumerar de forma exacta el soporte inducido por un programa y ha sido utilizado para expresar modelos probabilísticos de manera composicional en lenguajes funcionales \cite{ErwigKollmansberger2006,ScibiorGG2015}.

El ejemplo de la moneda y el dado presentado en la Introducción constituye el caso más simple de esta interpretación. Sus dos elecciones uniformes generan los doce pares posibles y cada par recibe probabilidad $1/12$. Más que repetir ese ejemplo, esta sección se concentra en la representación concreta de las distribuciones y en operaciones que permiten construir modelos con elecciones uniformes, pesos no uniformes y dependencias entre variables.

La implementación utilizada en esta memoria pertenece al paquete externo \texttt{probability}, publicado en Hackage. Por lo tanto, no forma parte de los módulos estándar de \textsf{Haskell} ni del paquete \texttt{base}. Su módulo público principal es \texttt{Numeric.Probability.Distribution}, importado habitualmente con el calificador \texttt{Dist}; este módulo proporciona el tipo \texttt{T}, sus instancias de \texttt{Functor}, \texttt{Applicative} y \texttt{Monad}, y constructores de distribuciones como \texttt{certainly}, \texttt{choose}, \texttt{uniform} y \texttt{fromFreqs} \cite{HackageProbabilityDistribution}.

La definición esencial del tipo \texttt{T} se muestra en el Código~\ref{lst:probability-distribution-representation}. Una distribución se almacena internamente como una lista de pares: el primer componente de cada par es un resultado posible y el segundo es su probabilidad asociada \cite{HackageProbabilityDistribution}.

\begin{lstlisting}[style=haskell,caption={Representación interna de una distribución dentro del paquete \texttt{probability}.},label={lst:probability-distribution-representation}]
newtype T prob a = Cons
    { decons :: [(a, prob)] }
\end{lstlisting}

El parámetro \texttt{a} corresponde al tipo de los resultados y \texttt{prob} al tipo usado para sus probabilidades. En particular, \texttt{Dist.T Rational a} representa resultados de tipo \texttt{a} con pesos racionales. El uso de \texttt{Rational} permite conservar exactamente fracciones como $1/3$ o $3/10$, sin introducir los errores de redondeo propios de una representación de punto flotante. Por ejemplo, el constructor \texttt{Dist.Cons} puede almacenar los resultados \texttt{x} e \texttt{y} en una lista con los pesos $1/4$ y $3/4$, respectivamente.

Esta lista registra ramas de cómputo y no necesariamente constituye una representación canónica de la distribución. Un mismo resultado puede aparecer varias veces si fue alcanzado mediante caminos distintos. Además, el orden de los pares refleja el orden en que esas ramas fueron generadas. La semántica probabilística se obtiene al considerar, para cada resultado, la suma de los pesos de todas las apariciones que lo producen.

Las operaciones del módulo construyen estas listas de manera controlada. Una elección determinista puede escribirse como \texttt{Dist.certainly x} y produce el único par $(\mathtt{x},1)$. La expresión \texttt{Dist.choose p x y} produce dos ramas: \texttt{x} con probabilidad $p$ e \texttt{y} con probabilidad $1-p$. Por su parte, \texttt{Dist.uniform} asigna peso $1/n$ a cada elemento de una lista de largo $n$, mientras que \texttt{Dist.fromFreqs} recibe pares de valores y frecuencias, y divide cada frecuencia por la suma total para obtener probabilidades. Así, las frecuencias $5$, $3$ y $2$ se transforman respectivamente en las probabilidades $1/2$, $3/10$ y $1/5$ \cite{HackageProbabilityDistribution}.

La instancia de \texttt{Monad} determina cómo se propagan estos pesos. Si una distribución contiene una rama $(x,p)$ y el cómputo posterior aplicado a \texttt{x} contiene una rama $(y,q)$, el operador \verb|(>>=)| incorpora el resultado \texttt{y} con peso $p q$. En consecuencia, la composición enumera las combinaciones posibles y multiplica las probabilidades a lo largo de cada camino. Cuando varios caminos terminan en el mismo valor, sus pesos permanecen inicialmente en pares separados.

El Código~\ref{lst:weighted-probability-example} combina las tres formas de construcción anteriores. Primero define una distribución ponderada para el clima. Luego usa \texttt{Dist.choose} para seleccionar una actividad según el clima obtenido y, finalmente, elige uniformemente su duración. La distribución retornada omite el clima, que actúa como una variable intermedia del modelo.

\noindent\begin{minipage}{\textwidth}
\begin{lstlisting}[style=haskell,caption={Modelo probabilístico con frecuencias, elecciones condicionadas y una distribución uniforme.},label={lst:weighted-probability-example}]
import qualified Numeric.Probability.Distribution as Dist

data Weather = Sunny | Cloudy | Rainy
    deriving (Eq, Ord, Show)

data Activity = Outdoor | Indoor
    deriving (Eq, Ord, Show)

weekendPlan :: Dist.T Rational (Activity, Int)
weekendPlan = do
    weather <- Dist.fromFreqs
        [(Sunny, 5), (Cloudy, 3), (Rainy, 2)]
    activity <- case weather of
        Sunny  -> Dist.choose (4 / 5) Outdoor Indoor
        Cloudy -> Dist.choose (1 / 2) Outdoor Indoor
        Rainy  -> Dist.choose (1 / 10) Outdoor Indoor
    hours <- Dist.uniform [1, 2, 3]
    return (activity, hours)

normalizedWeekendPlan = Dist.norm weekendPlan
\end{lstlisting}
\end{minipage}

Antes de normalizar, este programa genera dieciocho ramas: tres posibilidades para el clima, dos para la actividad y tres para la duración. Sin embargo, como el clima no forma parte del resultado, solo existen seis pares distintos de actividad y duración. Por ejemplo, el resultado \texttt{(Outdoor, 1)} puede alcanzarse desde cualquiera de los tres climas. La suma exacta de sus ramas es

\begin{ReportExpression}{Cálculo de la probabilidad del resultado \texttt{(Outdoor, 1)}.}{expr:probabilidad-outdoor-una-hora}
\[
\frac{1}{3}\left(
    \frac{1}{2}\frac{4}{5}
    + \frac{3}{10}\frac{1}{2}
    + \frac{1}{5}\frac{1}{10}
\right)
= \frac{19}{100}.
\]
\end{ReportExpression}

La función \texttt{Dist.norm} agrupa los resultados iguales y suma sus probabilidades, produciendo en este caso una lista canónica de seis resultados. En este módulo, esta operación de normalización corresponde específicamente al agrupamiento de resultados repetidos; no vuelve a dividir pesos arbitrarios para hacer que sumen uno. Esa segunda operación sí ocurre al construir una distribución mediante \texttt{Dist.fromFreqs} \cite{HackageProbabilityDistribution}.

La distinción entre distribución y lista interna es relevante para la conmutatividad. Al intercambiar dos cómputos probabilísticos independientes, el producto de sus probabilidades se conserva, aunque el recorrido de las ramas puede producir los pares en un orden distinto. Por ello, la equivalencia considerada en esta memoria se establece sobre las distribuciones normalizadas y no sobre el orden literal expuesto por \texttt{decons}. Bajo esta interpretación, la mónada de distribuciones discretas finitas constituye el caso conmutativo descrito en la sección anterior \cite[sección~3.1]{Jacobs2018}.

El mismo tipo \texttt{Dist.T Rational} se utiliza en el modelo académico del Código~\ref{lst:main-example}. Allí, \texttt{Dist.uniform} representa la preparación y la dificultad, mientras que \texttt{Dist.choose} expresa las probabilidades condicionadas de aprobar. Tanto ese ejemplo como el corpus probabilístico aplican \texttt{Dist.norm} antes de observar el resultado. De esta forma, la validación compara probabilidades racionales exactas y evita confundir diferencias en el orden interno o en la multiplicidad de las ramas con diferencias semánticas entre distribuciones.

% !TEX root = ../../main.tex

\section{ApplicativeDo}

\texttt{ApplicativeDo} es una extensión de \textsf{GHC} basada en el trabajo de Marlow, Peyton Jones, Kmett y Mokhov \cite{MarlowPJKM2016}. Su objetivo es permitir que programas escritos con \textit{do-notation} usen combinadores aplicativos cuando las dependencias lo permiten. De esta forma, el programador conserva la sintaxis monádica habitual, mientras el compilador puede generar una estructura de ejecución más paralelizable.

La transformación que implementa \texttt{ApplicativeDo} recibe un bloque \texttt{do}, analiza las dependencias entre sus statements y produce una expresión equivalente que utiliza operaciones aplicativas donde existe independencia, sin dejar de utilizar operaciones monádicas donde un cómputo depende del resultado de otro. Para ello, el algoritmo se organiza conceptualmente como un \textit{pipeline} de dos etapas: \textit{rearrange} construye un plan que combina secuencias y agrupaciones aplicativas, mientras que \textit{desugaring} traduce ese plan a combinadores explícitos \cite[pp.~95--96]{MarlowPJKM2016}.

Esta separación también forma parte del diseño de la implementación. La primera etapa conserva la estructura de los statements y registra cómo deben agruparse; la segunda consume esa información para producir las operaciones que implementan el plan. A continuación se profundizará en el diseño conceptual de las etapas ya mencionadas \textit{rearrange} y \textit{desugaring}, aplicando luego el algoritmo de \texttt{ApplicativeDo} al programa probabilístico del Código~\ref{lst:main-example}.

\subsection{Etapa de Rearrange}
\label{sec:applicative-do-rearrange}

Para presentar el algoritmo, considérese una expresión de \textit{do-notation} cuyos statements de enlace tienen la forma \verb|p_i <- e_i|. Cada statement se denota mediante $s_i$, donde $i$ indica su posición dentro del bloque. El algoritmo representa entonces la expresión mediante la forma abstracta \verb|do l e|, en la que $l=\{s_1;\ldots;s_n\}$ es la secuencia de statements y $e$ es la expresión final que produce el resultado. Esta última corresponde al argumento de un \texttt{return} o \texttt{pure} terminal y no participa en los cálculos de \textit{rearrange}.

La Figura~\ref{fig:rearrangement-marlow} presenta las funciones que componen esta etapa. La función principal, \texttt{rearrange}, recibe la secuencia $l$ y construye otra secuencia $l'$ que puede contener la forma paralela $(l_1\mid\cdots\mid l_k)$. Esta forma no pertenece a la sintaxis escrita por el programador, sino a una sintaxis abstracta intermedia que expresa que las secuencias $l_1,\ldots,l_k$ son independientes y pueden combinarse mediante operaciones aplicativas.

\begin{figure}[ht]
  \centering
  \begingroup
  \small
  \setlength{\jot}{2pt}
  \newcommand{\AdoFn}[1]{\mathsf{#1}}
  \newcommand{\AdoAppend}{\mathbin{\text{\ttfamily ++}}}
  \newcommand{\AdoLine}[2]{%
    \noindent\hspace*{#1}\(#2\)\par
  }
  \newcommand{\AdoBlockGap}{\vspace{1.5\baselineskip}}
  \begin{minipage}{0.6\textwidth}
  \setlength{\parindent}{0pt}
  \setlength{\parskip}{0pt}
  \raggedright
  \AdoLine{0pt}{
    \begin{alignedat}{2}
    \AdoFn{rearrange}\,\{s_1;\ldots;s_n\}
      &={} \{s_1\},
      &\qquad& \textit{if } n=1,\\
      &={} \AdoFn{split}\,g_1,
      && \textit{if } k=1,\\
      &={} \{(\AdoFn{split}\,g_1\mid\ldots\mid\AdoFn{split}\,g_k)\},
      && \text{otherwise}
    \end{alignedat}
  }
  \AdoLine{1.2em}{\text{\bfseries where}}
  \AdoLine{2.8em}{
    g_1\ldots g_k
      = \AdoFn{segments}\,\{s_1;\ldots;s_n\}
  }

  \AdoBlockGap
  \AdoLine{0pt}{
    \AdoFn{segments}\,\{s_1;\ldots;s_n\}
      = \{s_1;\ldots;s_{i_1}\}\ldots
        \{s_{(i_k)+1};\ldots;s_n\}
  }
  \AdoLine{1.2em}{\text{\bfseries where}}
  \AdoLine{2.8em}{
    \begin{aligned}
    i_1\ldots i_k
      &= \bigl\{\,i\in 1\ldots n\\
      &\quad \mathrel{|}
        \AdoFn{bv}\,\{s_1;\ldots;s_i\}
        \cap
        \AdoFn{fv}\,\{s_{i+1};\ldots;s_n\}
        =\emptyset\,\bigr\}
    \end{aligned}
  }

  \AdoBlockGap
  \AdoLine{0pt}{
    \begin{alignedat}{2}
    \AdoFn{split}\,\{s_1;\ldots;s_n\}
      &={} \{s_1\},
      &\qquad& \textit{if } n=1,\\
      &={} \AdoFn{splitat}\,i_{\mathrm{opt}},
      && \text{otherwise}
    \end{alignedat}
  }
  \AdoLine{1.2em}{\text{\bfseries where}}
  \AdoLine{2.8em}{
    \AdoFn{splitat}\,i
      = \AdoFn{rearrange}\,\{s_1;\ldots;s_i\}
        \AdoAppend
        \AdoFn{rearrange}\,\{s_{i+1};\ldots;s_n\}
  }
  \AdoLine{2.8em}{
    i_{\mathrm{opt}}\in 1\ldots n\ \text{\bfseries such that}
  }
  \AdoLine{3.6em}{
    \forall j.\ 1\leq j<n.\quad
    \AdoFn{cost}_s\,(\AdoFn{splitat}\,j)
      \geq \AdoFn{cost}_s\,(\AdoFn{splitat}\,i_{\mathrm{opt}})
  }

  \AdoBlockGap
  \AdoLine{0pt}{
    \AdoFn{cost}_s\,\{s_1;\ldots;s_n\}
      = \textstyle\sum\bigl\{\AdoFn{cost}_a\,s_i
        \mathrel{|} 1\leq i\leq n\bigr\}
  }

  \AdoBlockGap
  \AdoLine{0pt}{
    \begin{array}{@{}l@{\;}c@{\;}l@{}}
    \AdoFn{cost}_a\,(p\leftarrow e)
      & = & 1\\
    \AdoFn{cost}_a\,(l_1\mid\ldots\mid l_n)
      & = & \AdoFn{max}\bigl\{\AdoFn{cost}_s\,l_i
        \mathrel{|} 1\leq i\leq n\bigr\}
    \end{array}
  }

  \AdoBlockGap
  \AdoLine{0pt}{
    \begin{array}{@{}l@{\;}c@{\;}l@{}}
    \AdoFn{fv}\,\{s_1;\ldots;s_n\}
      & = & \text{\itshape the free variables of }s_1\ldots s_n\\
    \AdoFn{bv}\,\{s_1;\ldots;s_n\}
      & = & \text{\itshape the variables bound by }s_1\ldots s_n
    \end{array}
  }
  \end{minipage}
  \endgroup
  \captionsetup{skip=14pt}
  \caption{Etapa de \textit{rearrange} propuesta por Marlow et al. \cite[p.~95]{MarlowPJKM2016}.}
  \label{fig:rearrangement-marlow}
\end{figure}

Si la entrada contiene un único statement, \texttt{rearrange} lo retorna sin modificaciones. Para una secuencia de mayor tamaño, primero invoca a \texttt{segments}, que examina las $n-1$ separaciones posibles entre statements consecutivos. Una separación ubicada después de $s_i$ divide la entrada en el prefijo $\{s_1;\ldots;s_i\}$ y el sufijo $\{s_{i+1};\ldots;s_n\}$. El corte es válido cuando se cumple:

\begin{ReportExpression}{Condición de independencia entre el prefijo y el sufijo de una secuencia.}{expr:condicion-independencia-corte}
\[
\operatorname{bv}(\{s_1;\ldots;s_i\})
\cap
\operatorname{fv}(\{s_{i+1};\ldots;s_n\})
=\emptyset,
\]
\end{ReportExpression}

En la notación $\operatorname{bv}$ representa las variables vinculadas por una secuencia y $\operatorname{fv}$ sus variables libres. Si la intersección es vacía, ninguna variable producida por el prefijo es requerida por el sufijo y, por tanto, no existe una dependencia que atraviese esa separación. Al considerar todos los cortes válidos, \texttt{segments} divide la entrada en los grupos $g_1,\ldots,g_k$ que ya pueden formar las ramas de una agrupación aplicativa.

La condición anterior expresa una relación \textit{def-use}. Esta relación establece que el uso de una variable debe encontrarse después de su definición y dentro del alcance introducido por ella. Si una variable definida en $\{s_1;\ldots;s_i\}$ aparece libre en $\{s_{i+1};\ldots;s_n\}$, ambos fragmentos no pueden tratarse como ramas independientes: el segundo necesita un resultado producido por el primero. En cambio, una intersección vacía indica que el corte no separa ninguna definición de sus usos. Como la secuencia original ya está bien formada, una variable definida en el sufijo tampoco puede ser utilizada por un statement anterior del prefijo.

Cuando \texttt{segments} obtiene más de un grupo, \texttt{rearrange} construye la forma $(\texttt{split}\ g_1\mid\cdots\mid\texttt{split}\ g_k)$. Cada grupo se procesa mediante \texttt{split} para descubrir posibles agrupaciones anidadas. Si, por el contrario, toda la entrada constituye un único segmento, no existe un corte exterior que permita combinar dos partes aplicativamente y \texttt{rearrange} delega directamente la secuencia completa a \texttt{split}.

En su formulación óptima, la función \texttt{split} considera nuevamente las separaciones $i\in\{1,\ldots,n-1\}$, pero con un propósito distinto. Para cada $i$, aplica \texttt{rearrange} recursivamente al prefijo y al sufijo, y luego compone ambos resultados de manera secuencial. Entre estas alternativas selecciona aquella que minimiza el costo definido por el modelo estructural de la Expresión~\ref{expr:modelo-costos-estructural}. De esta manera, \texttt{segments} prioriza la independencia exterior que puede expresarse directamente mediante una agrupación aplicativa, mientras que \texttt{split} determina dónde conservar una composición monádica cuando una dependencia impide realizar esa agrupación.

Para realizar esta tarea, \textsf{GHC} implementa dos estrategias. La estrategia predeterminada evita comparar exhaustivamente todos los puntos de separación de cada segmento indivisible y selecciona uno mediante una heurística, por lo que presenta una complejidad temporal de $O(n^2)$. Aunque esta estrategia suele encontrar el plan de menor costo, no lo garantiza en todos los casos. Por otra parte, la opción de compilación \texttt{-foptimal-applicative-do} activa una estrategia basada en programación dinámica, que reutiliza los resultados calculados para las subsecuencias y garantiza la obtención de un plan de costo estructural mínimo. Esta variante tiene complejidad temporal $O(n^3)$ \cite{GHCApplicativeDo}.

La selección entre ambas estrategias se realiza antes de comenzar la construcción del plan de ejecución. Las dos comparten la lógica de \texttt{segments} para identificar grupos independientes, pero difieren en la estrategia utilizada para dividir los segmentos indivisibles mediante \texttt{split}. Si $n$ es la cantidad de statements procesados por \textit{rearrange}, sus complejidades temporales se resumen como:

\begin{ReportExpression}{Complejidad temporal de las variantes de \textit{rearrange}.}{expr:complejidad-rearrange}
\[
T_{\mathrm{rearrange}}(n)=
\begin{cases}
O(n^2), & \text{variante heurística predeterminada},\\
O(n^3), & \text{variante óptima con programación dinámica}.
\end{cases}
\]
\end{ReportExpression}

Finalmente, el resultado de la etapa \textit{rearrange} es la expresión abstracta \verb|do l' e|. Aunque $l'$ puede tener una estructura de árbol con secuencias y agrupaciones aplicativas, el algoritmo no permuta statements: si se eliminan las formas paralelas y se concatenan sus ramas de izquierda a derecha, se recupera la secuencia original $\{s_1;\ldots;s_n\}$ \cite[p.~95]{MarlowPJKM2016}.

\subsection{Etapa de Desugaring}
\label{sec:applicative-do-desugaring}

La segunda etapa recibe el plan $l'$ construido por \textit{rearrange} junto con la expresión final $e$. Su función principal puede representarse mediante la firma:

\begin{ReportExpression}{Firma abstracta de la función \texttt{desugar}.}{expr:firma-desugar}
\verb|desugar :: Stmts -> Expr -> Expr|
\end{ReportExpression}

El segundo argumento actúa como la continuación del plan: es la expresión que debe evaluarse una vez ejecutados los efectos de la secuencia y vinculadas las variables que utiliza.

La sintaxis intermedia permite que un statement sea un enlace ordinario o una agrupación $(l_1\mid\cdots\mid l_n)$. En esta última, cada $l_i$ es a su vez una secuencia de statements. Las variables definidas dentro de una rama no quedan disponibles para las demás ramas, pero aquellas requeridas por la continuación deben ser entregadas como resultados de esa rama. Esta restricción es la que permite reemplazar la agrupación por una expresión construida con \verb|(<$>)| y \verb|(<*>)|.

La Figura~\ref{fig:desugaring-marlow} define el proceso mediante cinco casos recursivos. En conjunto, estos casos eliminan la sintaxis intermedia y producen una expresión formada por operaciones de \texttt{Functor}, \texttt{Applicative} y \texttt{Monad}.

\begin{figure}[ht]
  \centering
  \begingroup
  \small
  \setlength{\jot}{2pt}
  \shorthandoff{<>}
  \newcommand{\AdoFn}[1]{\mathsf{#1}}
  \newcommand{\AdoFmap}{\mathbin{\text{\ttfamily <\$>}}}
  \newcommand{\AdoApply}{\mathbin{\text{\ttfamily <*>}}}
  \newcommand{\AdoBind}{\mathbin{\text{\ttfamily >>=}}}
  \newcommand{\AdoEquals}{\mathrel{\text{\ttfamily ==}}}
  \newcommand{\AdoLine}[2]{%
    \noindent\hspace*{#1}\(#2\)\par
  }
  \newcommand{\AdoTaggedLine}[2]{%
    \noindent\makebox[\linewidth][l]{\(#1\)\hfill\textit{(#2)}}\par
  }
  \newcommand{\AdoBlockGap}{\vspace{1.5\baselineskip}}
  \begin{minipage}{0.6\textwidth}
  \setlength{\parindent}{0pt}
  \setlength{\parskip}{0pt}
  \raggedright
  \AdoLine{0pt}{
    \AdoFn{desugar}
      :: \mathit{Stmts}\to\mathit{Expr}\to\mathit{Expr}
  }

  \AdoBlockGap
  \AdoTaggedLine{
    \AdoFn{desugar}\,\{\}\,e = \AdoFn{pure}\,e
  }{0}

  \AdoBlockGap
  \AdoTaggedLine{
    \AdoFn{desugar}\,\{p\leftarrow e\}\,e'
  }{1}
  \AdoLine{2.4em}{
    \begin{array}{@{}l@{\;}c@{\;}l@{}}
    \mid\ p\AdoEquals e'
      & = & e\\
    \mid\ \text{otherwise}
      & = & (\lambda p\to e')\AdoFmap e
    \end{array}
  }

  \AdoBlockGap
  \AdoTaggedLine{
    \AdoFn{desugar}\,\{p\leftarrow e;\,l\}\,e'
      = e\AdoBind\lambda p\to\AdoFn{desugar}\,l\,e'
  }{2}

  \AdoBlockGap
  \AdoTaggedLine{
    \AdoFn{desugar}\,\{(l_1\mid\ldots\mid l_n)\}\,e
  }{3}
  \AdoLine{2.4em}{
    = (\lambda p_1\ldots p_n\to e)
      \AdoFmap e_1\AdoApply\ldots\AdoApply e_n
  }
  \AdoLine{2.4em}{
    \text{\bfseries where}\; (p_i,e_i)
      = \AdoFn{desugar}_{\mathit{arg}}\,l_i\,\AdoFn{fv}(e)
  }

  \AdoBlockGap
  \AdoTaggedLine{
    \AdoFn{desugar}\,\{(l_1\mid\ldots\mid l_n);\,s\}\,e
  }{4}
  \AdoLine{2.4em}{
    = \AdoFn{join}\bigl((\lambda p_1\ldots p_n\to e')
      \AdoFmap e_1\AdoApply\ldots\AdoApply e_n\bigr)
  }
  \AdoLine{2.4em}{\text{\bfseries where}}
  \AdoLine{5em}{
    e' = \AdoFn{desugar}\,s\,e
  }
  \AdoLine{5em}{
    (p_i,e_i)
      = \AdoFn{desugar}_{\mathit{arg}}\,l_i\,\AdoFn{fv}(e')
  }

  \AdoBlockGap
  \AdoLine{0pt}{
    \AdoFn{desugar}_{\mathit{arg}}
      :: \mathit{Stmts}\to\mathit{Set}\;\mathit{Var}
        \to(\mathit{Pat},\mathit{Expr})
  }
  \AdoLine{0pt}{
    \AdoFn{desugar}_{\mathit{arg}}\,\{p\leftarrow e\}\,vs
      = (p,e)
  }
  \AdoLine{0pt}{
    \AdoFn{desugar}_{\mathit{arg}}\,l\,vs
      = \bigl((v_1,\ldots,v_k),
        \AdoFn{desugar}\,l\,(v_1,\ldots,v_k)\bigr)
  }
  \AdoLine{2.4em}{\text{\bfseries where}}
  \AdoLine{2.4em}{
    v_1\ldots v_k = \AdoFn{bv}(l)\cap vs
  }
  \end{minipage}
  \endgroup
  \captionsetup{skip=14pt}
  \caption{Etapa de \textit{desugaring} propuesta por Marlow et al. \cite[p.~96]{MarlowPJKM2016}.}
  \label{fig:desugaring-marlow}
\end{figure}

\begin{description}
  \item[Caso (0).] Si la secuencia de statements está vacía, no queda ningún efecto por ejecutar y la continuación se introduce en el contexto mediante \texttt{pure}.

  \item[Caso (1).] Si solo queda un enlace \verb|p <- e|, el resultado de $e$ se transforma mediante \verb|(<$>)| para producir la continuación. Cuando la continuación coincide exactamente con el patrón $p$, el algoritmo retorna $e$ directamente y evita introducir una transformación identidad innecesaria.

  \item[Caso (2).] Si un enlace \verb|p <- e| está seguido por más statements, estos pueden depender del valor asociado a $p$. Por ello se utiliza \verb|(>>=)| y el resto de la secuencia se desazucara recursivamente dentro de la función que recibe dicho valor.

  \item[Caso (3).] Si el único elemento restante es una agrupación $(l_1\mid\cdots\mid l_n)$, cada rama se convierte en un argumento independiente. Sus resultados se reciben mediante una función anónima que construye la continuación, y las expresiones de las ramas se combinan usando \verb|(<$>)| y una sucesión de \verb|(<*>)|.

  \item[Caso (4).] Si una agrupación aplicativa está seguida por otros statements, primero se desazucara ese resto para obtener una continuación que todavía produce un cómputo contextualizado. La combinación aplicativa genera entonces un cómputo anidado, por lo que \texttt{join} elimina un nivel del contexto y permite continuar con el plan.
\end{description}

La función auxiliar \texttt{desugar\_arg} determina el patrón y la expresión correspondientes a cada rama de una agrupación. Para ello, intersecta las variables definidas por la rama con las variables libres de la continuación. Así, una rama compuesta puede realizar varios enlaces internos, pero solo entrega al exterior los resultados que serán utilizados posteriormente. Luego, su propia secuencia se procesa recursivamente mediante \texttt{desugar}, lo que permite conservar dependencias monádicas dentro de una rama y combinarlas aplicativamente con las demás.

\subsection{Aplicación del algoritmo}
\label{sec:applicative-do-aplicacion}

Para ilustrar ambas etapas se retomará el modelo probabilístico del estudiante presentado en el Código~\ref{lst:main-example}. Con el fin de concentrar la derivación en la estructura del programa, las cuatro expresiones monádicas del lado derecho de los enlaces se representarán temporalmente mediante \texttt{expr1}, \texttt{expr2}, \texttt{expr3} y \texttt{expr4}:

\begin{lstlisting}[style=haskell,caption={Representación abreviada del modelo probabilístico para aplicar \texttt{ApplicativeDo}.},label={lst:applicative-do-main-example-abstract}]
    studentResults = do
        preparation <- expr1
        quiz <- expr2
        difficulty <- expr3
        finalExam <- expr4
        return (preparation, quiz, difficulty, finalExam)
\end{lstlisting}

Estas expresiones son metavariables y no definiciones independientes de \textsf{Haskell}: \texttt{expr2} conserva el uso libre de \texttt{preparation}, mientras que \texttt{expr4} conserva los usos de \texttt{preparation} y \texttt{difficulty}. Por su parte, \texttt{expr1} y \texttt{expr3} no utilizan variables definidas por otros statements del bloque. Denotando los cuatro statements por $s_1$, $s_2$, $s_3$ y $s_4$, respectivamente, la entrada abstracta de \textit{rearrange} es:

\begin{ReportExpression}{Entrada abstracta de \texttt{rearrange} para el modelo probabilístico del estudiante.}{expr:entrada-rearrange-ejemplo}
\[
\texttt{do}\ \{s_1;s_2;s_3;s_4\}\
(\texttt{preparation}, \texttt{quiz}, \texttt{difficulty}, \texttt{finalExam}).
\]
\end{ReportExpression}

La llamada inicial a \texttt{segments} examina los tres posibles puntos de corte. El corte después de $s_1$ no es válido porque \texttt{preparation} se utiliza en $s_2$ y $s_4$. El corte después de $s_2$ tampoco lo es, pues $s_4$ todavía utiliza \texttt{preparation}. Finalmente, el corte después de $s_3$ separaría las definiciones de \texttt{preparation} y \texttt{difficulty} de sus usos en $s_4$. En consecuencia, los cuatro statements forman un único segmento y se aplica \texttt{split}.

Para las separaciones $i\in\{1,2,3\}$, las llamadas recursivas a \texttt{rearrange} producen los siguientes planes y costos:

\begin{ReportExpression}{Planes y costos obtenidos por \texttt{rearrange} para las tres separaciones posibles.}{expr:resultados-rearrange-ejemplo}
\[
\begin{aligned}
i=1 &\quad\Longrightarrow\quad
  \{s_1;(s_2\mid\{s_3;s_4\})\}
  && C=3,\\
i=2 &\quad\Longrightarrow\quad
  \{s_1;s_2;s_3;s_4\}
  && C=4,\\
i=3 &\quad\Longrightarrow\quad
  \{(\{s_1;s_2\}\mid s_3);s_4\}
  && C=3.
\end{aligned}
\]
\end{ReportExpression}

La primera y la tercera alternativa tienen el mismo costo mínimo. Para continuar la derivación se toma la primera, que coincide con el plan producido por \textsf{GHC} para el orden original, mostrado también en la Tabla~\ref{tab:main-example-plans}. Al sustituir nuevamente los patrones de los statements, el objeto entregado a la etapa de \textit{desugaring} tiene la siguiente estructura intermedia:

\begin{lstlisting}[style=haskell,caption={Plan intermedio producido por la etapa de \textit{rearrange}.},label={lst:applicative-do-main-example-rearranged}]
    do {
        preparation <- expr1;
        (quiz <- expr2 |
            { difficulty <- expr3;
              finalExam <- expr4 })
    } (preparation, quiz, difficulty, finalExam)
\end{lstlisting}

El primer enlace permanece secuencial porque los dos componentes posteriores necesitan \texttt{preparation}. Después de obtener ese valor, la expresión asociada al quiz puede combinarse aplicativamente con la secuencia que obtiene la dificultad y, a partir de ella, el resultado del examen final. Al aplicar las reglas de \textit{desugaring}, se obtiene primero la expresión abreviada:

\newpage
\begin{lstlisting}[style=haskell,caption={Resultado del desazucarado expresado mediante \texttt{expr1}--\texttt{expr4}.},label={lst:applicative-do-main-example-desugared-abstract}]
    studentResults =
        expr1 >>= \preparation ->
            (\quiz (difficulty, finalExam) ->
                (preparation, quiz, difficulty, finalExam))
            <$> expr2
            <*> (expr3 >>= \difficulty ->
                    (\finalExam -> (difficulty, finalExam))
                    <$> expr4)
\end{lstlisting}

En esta expresión, \texttt{expr2} se evalúa dentro del alcance de \texttt{preparation}. La segunda rama conserva un enlace monádico entre \texttt{expr3} y \texttt{expr4}, ya que la última requiere \texttt{difficulty}; el par producido por esa rama permite entregar ambos valores a la función ubicada a la izquierda de \verb|(<$>)|. Este caso no requiere \texttt{join}, porque la agrupación aplicativa es el último componente del plan.

Finalmente, al reemplazar cada metavariable por la expresión correspondiente del programa original, se obtiene la forma explícita generada conceptualmente por las dos etapas de \texttt{ApplicativeDo}:

\begin{lstlisting}[style=haskell,caption={Desazucarado aplicativo del modelo probabilístico del estudiante.},label={lst:applicative-do-main-example-desugared}]
    studentResults =
        Dist.uniform [Low, High] >>= \preparation ->
            (\quiz (difficulty, finalExam) ->
                (preparation, quiz, difficulty, finalExam))
            <$> quizGiven preparation
            <*> (Dist.uniform [Easy, Hard] >>= \difficulty ->
                    (\finalExam -> (difficulty, finalExam))
                    <$> examGiven preparation difficulty)
\end{lstlisting}

La expresión resultante combina operaciones monádicas y aplicativas según las dependencias del programa, pero al aplanar su plan conserva el orden relativo $[1,2,3,4]$. La posibilidad de evaluar reordenamientos válidos bajo una hipótesis explícita de conmutatividad queda fuera del algoritmo descrito y constituye la limitación que se desarrolla en el capítulo siguiente.

